\documentclass[runningheads]{llncs}

\usepackage[T1]{fontenc}
\usepackage{newtxtext}
\usepackage[varvw]{newtxmath}
\usepackage{graphicx}

\begin{document}

\title{Leer la curva de convergencia: dynamic time warping para el perfil de GWO--DE}
\titlerunning{Perfil de GWO--DE}

\author{Lucas Erazo\inst{1,2} \and Emanuel Vega\inst{1} \and Ricardo Soto\inst{1}}
\authorrunning{Erazo et al.}

\institute{Escuela de Ingeniería Informática, Pontificia Universidad Católica de Valparaíso, Valparaíso, Chile\\
\email{lucas.erazo@pucv.cl, emanuel.vega@pucv.cl, ricardo.soto@pucv.cl}\\
\and
Escuela de Ingeniería de Construcción y Transporte, Pontificia Universidad Católica de Valparaíso, Valparaíso, Chile}

\maketitle

\begin{abstract}
La historia del mejor valor encontrado se lee mientras la búsqueda sigue.
El cambio de perfil de GWO--DE sale de esa lectura, no de un contador fijo ni de una política entrenada.
\keywords{Dynamic time warping \and GWO--DE \and Perfil de búsqueda \and Mochila multidimensional.}
\end{abstract}

\section{Introducción}
Durante la búsqueda, la curva de convergencia se usa como señal mientras el proceso sigue, no solo como gráfico final.

\section{Trabajo relacionado}
% PENDIENTE: trabajo relacionado

\section{Lectura de la curva}
% PENDIENTE - Figura: ventana reciente contra la rampa y la constante. Ubicación: sección Lectura de la curva. Propósito: mostrar las dos referencias ancladas en el primer valor de la ventana. Fuentes: definición ya escrita en esta sección.

\subsection{Serie y ventana}
Sea $q_t$ el mejor valor registrado en la iteración $t$.
Para comparar maximización y minimización con el mismo sentido de progreso, la serie monitoreada es $x_t = q_t$ en el primer caso y $x_t = -q_t$ en el segundo.
Desde la iteración $W = 40$, el monitor no lee la corrida entera: toma la ventana reciente
\[
\mathbf{X} = [x_{t-W+1}, \ldots, x_t].
\]

\subsection{Referencias y distancias}
Esa ventana se compara con dos referencias del mismo largo, construidas para ella y ancladas en su primer valor.
No salen de corridas anteriores.
La de estancamiento es constante.
La rampa crece desde el mismo punto con pendiente fija $s_{\min} = 2$, en unidades de la serie por iteración.
Esa pendiente dibuja la referencia.
No obliga al algoritmo a mejorar al menos eso.
Si se deja en cero, pasa a ser el uno por ciento del rango de la ventana.
La Tabla~\ref{tab:monitor} fija el resto de la configuración de este híbrido.

\begin{table}
\caption{Configuración del monitor en GWO--DE.}
\label{tab:monitor}
\begin{tabular}{ll}
\hline
Parámetro & Valor \\
\hline
Ventana $W$ & 40 \\
Banda de alineamiento $w$ & 4 \\
Iteraciones sin mejora, $K_{\max}$ & 15 \\
Paciencia de confirmación $P$ & 8 \\
Pendiente de la rampa $s_{\min}$ & 2 \\
Percentil bajo / alto & 40 / 60 \\
Distancia por defecto & DDTW \\
\hline
\end{tabular}
\end{table}

Esos enteros no hacen el mismo trabajo.
$W = 40$ decide cuánta historia reciente entra en la comparación.
$K_{\max} = 15$ exige quince iteraciones seguidas sin mejora estricta antes de que las distancias puedan avanzar la confirmación.
$P = 8$ exige que la condición completa se sostenga ocho pasos, de modo que una distancia favorable en una sola iteración no cambia el perfil.
En este híbrido esa paciencia vale para toda la corrida.

La distancia entre la ventana y cada referencia se calcula con dynamic time warping, que tolera que la mejora no caiga en el mismo instante~\cite{berndt1994using}.
El alineamiento queda dentro de una banda de Sakoe--Chiba de ancho $w = \lfloor 0{,}1 \, W \rfloor$, aquí igual a 4, para que una deformación extrema no cuente como semejanza.
% PENDIENTE: fuente Sakoe--Chiba
La configuración de la Tabla~\ref{tab:monitor} usa la versión derivada, que aplica esa distancia a las diferencias de primer orden y hace que un desplazamiento vertical pese menos que el cambio de forma~\cite{keogh2001derivative}.
$D_1$ es la distancia a la rampa, $D_2$ la distancia a la constante y $\Delta = D_1 - D_2$.
La misma ventana puede leerse con DTW.
Este texto no compara las dos.

Una ventana parecida a la rampa tiende a quedar más cerca de esa referencia que de la constante, y una ventana casi plana tiende a quedar cerca de la constante.
Son tendencias, no una regla de cambio: el signo de $\Delta$ depende de la pendiente, del largo de la ventana y de la trayectoria observada.

\subsection{Decisión y confirmación}
Los umbrales no son cortes fijos.
Después de diez observaciones, el percentil 40 de $D_2$ marca cercanía a la constante, y el percentil 60 de $D_1$ y de $\Delta$ marca lejanía respecto de la rampa.
Con $N$ el tramo seguido sin mejora estricta, y con $K_{\max}$ tomado de la Tabla~\ref{tab:monitor}, la condición instantánea es
\begin{equation}
\mathrm{Estancado}(t) \iff
(N \ge K_{\max})
\land (D_2 \le \theta_c)
\land \bigl( (D_1 \ge \theta_r) \lor (\Delta \ge \theta_{\Delta}) \bigr).
\label{eq:estancado}
\end{equation}
La ecuación~(\ref{eq:estancado}) no autoriza el cambio.
Tiene que sostenerse $P$ pasos.
Si falla, la racha vuelve a cero.

\section{Cambio de perfil}
\subsection{Dos perfiles}
El monitor no elige otra metaheurística: cambia el perfil de GWO--DE.
Hay dos perfiles fijos.
Explotación usa $a = 0{,}5$, $F = 0{,}5$ y $CR = 0{,}9$.
Exploración usa $a = 2$, $F = 0{,}9$ y $CR = 0{,}3$.
La búsqueda empieza en explotación.

\subsection{Ida, vuelta y reinicio}
Si la explotación está activa y la ecuación~(\ref{eq:estancado}) se sostiene $P$ pasos, el paso siguiente usa exploración.
Si ya está en exploración y una observación lista muestra un tramo sin mejora menor que $K_{\max}$, vuelve a explotación.
Si sigue estancado, permanece ahí.
Después de cada cambio el monitor se reemplaza: se descartan la ventana, la racha y los umbrales, y no hay otra decisión hasta reunir $W$ observaciones nuevas.
La población y el mejor global no se reinician.
El cambio decidido en la generación $g$ rige desde $g+1$.
El mismo híbrido tiene otra regla binaria, que entra a exploración solo si $D_2 \le \theta_c$.
No es la de este texto.

\section{Contraste en la mochila}
\subsection{Instancias y corridas}
El contraste usa nueve instancias de la mochila multidimensional de Chu y Beasley~\cite{chu1998genetic}, con 100, 250 o 500 ítems y 5, 10 o 30 restricciones, cada una con un valor de referencia conocido.
Son 31 corridas independientes por condición, con la configuración de la Tabla~\ref{tab:monitor}.

\subsection{Controles}
La Tabla~\ref{tab:controles} separa el perfil adaptativo de dos controles que no leen la curva y no cambian: explotación fija y exploración fija.
El adaptativo solo deja la explotación si la ecuación~(\ref{eq:estancado}) se sostiene $P$ pasos, y solo vuelve si la exploración muestra mejora reciente.

\begin{table}
\caption{Qué observa cada perfil de GWO--DE y cuándo cambia.}
\label{tab:controles}
\begin{tabular}{lll}
\hline
Procedimiento & Qué observa & Cuándo cambia \\
\hline
Perfil adaptativo & Forma de la ventana y persistencia & Disparo completo \\
Explotación fija & Nada de la curva & No cambia \\
Exploración fija & Nada de la curva & No cambia \\
\hline
\end{tabular}
\end{table}

Esas dos diferencias aíslan la pregunta.
Si el perfil adaptativo no supera a la explotación fija, salir de ese modo no aporta.
Si no supera a la exploración fija, el regreso tampoco.
Una diferencia solo cuenta si aparece en la calidad final y no se explica por haber usado un perfil distinto desde el inicio.
Si la misma lectura se sostiene en otro dominio no se examina aquí.

\section{Resultados}
% PENDIENTE: contraste numérico en la mochila. No inventar promedios, brechas ni pruebas.

\section{Discusión}
% PENDIENTE

\section{Conclusión}
% PENDIENTE

\begin{credits}
\subsubsection{\ackname}
% PENDIENTE
\subsubsection{\discintname}
% PENDIENTE
\end{credits}

\begin{thebibliography}{3}
\bibitem{berndt1994using}
Berndt, D.J., Clifford, J.: Using dynamic time warping to find patterns in time series. In: KDD Workshop (AAAI-94 Workshops on Knowledge Discovery in Databases), vol. 398, pp. 359--370 (1994)

\bibitem{keogh2001derivative}
Keogh, E.J., Pazzani, M.J.: Derivative dynamic time warping. In: Proceedings of the 2001 SIAM International Conference on Data Mining, pp. 1--11 (2001)

\bibitem{chu1998genetic}
Chu, P.C., Beasley, J.E.: A genetic algorithm for the multidimensional knapsack problem. Journal of Heuristics 4(1), 63--86 (1998)
\end{thebibliography}

\end{document}
